\medskip

L'espace est muni d'un repère orthonormé \Oijk. On considère :

\begin{itemize}[label=$\bullet~$]
\item les points A$(-1~;~2~;~1)$, \:B$(1~;~-1~;~2)$ et C(1~;~1~;~1) ;
\item la droite $d$ dont une représentation paramétrique est donnée par:

\[d : \:\: \left\{\begin{array}{l c l}
x& =& \frac 32 + 2t\\
y&=&2 +t\\
z&=&3 - t 
\end{array}\right. \:\text{avec}\: t \in\R ;\]
\item la droite $d'$ dont une représentation paramétrique est donnée par :

\[d : \:\: \left\{\begin{array}{l c l}
x& =&\phantom{3 - 2}s\\
y&=&\frac 32 + s\\
z&=&3 - 2s 
\end{array}\right. \:\text{avec}\: s \in\R ;\]
\end{itemize}

\medskip

\begin{center}\textbf{Partie A} \end{center}

\smallskip

\begin{enumerate}
\item Montrer que les droites $d$ et $d'$ sont sécantes au point S$\left(- \frac 12~;~1~;~4\right)$.
\item 
	\begin{enumerate}
		\item Montrer que le vecteur $\vect{n}\begin{pmatrix} 1\\2\\4\end{pmatrix}$ est un vecteur normal au plan (ABC) 
		\item En déduire qu'une équation cartésienne du plan (ABC) est:

\[x + 2y + 4z - 7 = 0.\]
	\end{enumerate}
\item Démontrer que les points A, B, C et S ne sont pas coplanaires.
\item 
	\begin{enumerate}
		\item Démontrer que le point H$(-1~;~0~;~2)$ est le projeté orthogonal de S sur le plan (ABC)
		\item En déduire qu'il n'existe aucun point $M$ du plan (ABC) tel que S$M < \dfrac{\sqrt{21}}{2}$.
	\end{enumerate}
\end{enumerate}

\smallskip

\begin{center}\textbf{Partie B} \end{center}

On considère un point $M$ appartenant au segment [CS]. On a donc $\vect{\text{C}M} = k\vect{\text{CS}}$ avec  $k$ réel de l'intervalle [0~;~1].

\medskip

\begin{enumerate}
\item Déterminer les coordonnées du point $M$ en fonction de $k$.
\item Existe-t-il un point $M$ sur le segment [CS] tel que le triangle ($M$AB) soit rectangle
en $M$ ?
\end{enumerate}

\bigskip
